% !PRINTED 6
\weberpage{6}
%<<CONT
is not continuous. For take any rational fraction $\mu=m:n$, in which
$m$ and $n$ have no common divisor and are not both square numbers.
Then $\mu$ is not the square of a rational fraction. For if
$\mu=p^2:q^2$, then $mq^2=np^2$ would follow, and from this $m=p^2$,
$n=q^2$, contrary to the assumption. Therefore if we form a cut
$(A,B)$ in the manifold $\mathfrak R$ by assigning to $A$ every
element $a$ of $\mathfrak R$ whose square is smaller than $\mu$, and
to $B$ every element $b$ whose square is greater than $\mu$, then
neither $A$ contains a greatest element nor $B$ a least element, and
the cut $(A,B)$ is not produced by an element of $\mathfrak R$.

For suppose that $a=p:q$ is any element of $A$, so that
$p^2:q^2<m:n$, or $np^2<mq^2$. Let $y$ be any natural number and
choose another natural number $x$ so that $x>y$ and
$x(mq^2-np^2)>ny(2p+1)$. Then
\[
 x^2(mq^2-np^2)>nxy(2p+1)>n(2pxy+y^2),
\]
and hence also
\[
 mq^2x^2>n(px+y)^2.
\]
Thus, if we put $a'=(px+y):qx$, then $a'>a$ and $a'^2<\mu$;
hence $a'$ too is contained in $A$. In exactly the same way one can
show that there is no least element in $B$.

The manifold $\mathfrak R$ can, however, serve as our starting point
for constructing a continuous set. The totality of all cuts in
$\mathfrak R$ is certainly a manifold, which may be denoted by
$\mathfrak S$. Consider two different elements of it,
$\alpha=(A,B)$ and $\alpha'=(A',B')$. Then either $A$ is a part of
$A'$ or $A'$ is a part of $A$. For if some

\par\vfill
\noindent\rule{0.22\linewidth}{0.4pt}

{\footnotesize\noindent
manifolds to which continuity belongs were formed by
\gesperrt{Weierstrass} and G.\ \gesperrt{Cantor}.

The definition of continuity, as we here lay it down after
\gesperrt{Dedekind}, is exhaustive in the following respect: a set
continuous in this sense, if it also possesses the property of
measurability presently to be discussed, cannot be a part of a richer
continuous set. I do not know whether this property has already been
proved anywhere, and hope to return to it on another occasion. I
remark, however, that such a property can be proved only for measurable
sets. A merely ordered set can always, however dense it may be, be
regarded as a part of one still denser.}
%CONT>>
